\chapter{透镜、2f 与 4f 系统}\label{ch:lens}

\begin{idea}[一句话版本]
\textbf{透镜 = 乘一个负二次相位}；\textbf{透镜前焦面 $\to$ 后焦面 = 一次傅里叶变换}；
\textbf{两次傅里叶变换 = 倒像}，所以 4f 系统就是“变换—在频谱面动手脚—再变换回来”。
做 4f 类题只需三步：写频谱 $\to$ 乘滤波器/相位 $\to$ 再变换并代入 $\F\{\F\{f\}\}=f(-x)$。
\end{idea}

\begin{idea}[本章题型一览（先找到自己的题，再跳到对应小节）]
\begin{tabular}{@{}P{6.2cm}P{5.6cm}P{3.4cm}@{}}
题目长什么样 & 本章单元 & 往年\\ \hline
给一个相位/厚度分布，问等效于什么透镜、焦距多少 & 等效透镜（§\ref{sec:equivlens}）& 2025-3(i)\\
物在透镜前/后某处，求焦面复振幅、强度、谱点位置 & 焦面场（§\ref{sec:focalplane}）& \\
4f 求输出、放大率；级联 & 4f 输出（§\ref{sec:4fderiv}）& 2024-4(i)\\
频谱面放滤波器，求 CTF / 输出 & 滤波与 CTF（§\ref{sec:4fctf}）& 2024-4(ii)\\
输入光栅，频谱面挡掉某些级，求输出条纹 & 光栅滤波（§\ref{sec:gratingfilter}）& \\
频谱面加线性相位（棱镜、SLM 倾斜相位）& 输出平移（§\ref{sec:linphase}）& 2025-3(ii)\\
输入面偏离前焦面，证明 $x_2=M^2x_1$ & 纵向偏离（§\ref{sec:4foffset}）& 2024-5\\
Alvarez 透镜平移变焦；4f 实现 & Alvarez（§\ref{sec:alvarez}）& 2025-3\\
光学加减法、相衬、暗场、高通/低通 & 相干处理（§\ref{sec:processing}）& \\
\end{tabular}
\end{idea}

\needspace{8\baselineskip}\section{透镜的相位变换}\label{sec:lensphase}

\begin{formula}{薄透镜复振幅透过率}\ix{透镜!相位因子}\ix{透镜!透过率}
\[
t_L(x,y)=P(x,y)\exp\!\left(-ik\frac{x^2+y^2}{2f}\right),\qquad k=\frac{2\pi}{\lambda}
\]
$P(x,y)$：光瞳函数（孔内为 1，孔外为 0；有像差时 $P=|P|\ee^{ikW}$）。负号 = 会聚；$f<0$ 时为发散透镜。

来源（课件 [F1-p16]）：球面近似为抛物面 $z=\frac{x^2+y^2}{2R}$；$(x,y)$ 处比中心少走玻璃 $\frac{x^2+y^2}{2R_1}+\frac{x^2+y^2}{2R_2}$，少的光程 $(n-1)(\cdots)$，
得 $t=\exp\!\left[-ik(n-1)\left(\frac1{R_1}+\frac1{R_2}\right)\frac{x^2+y^2}{2}\right]$，
即\textbf{磨镜者公式（课件写法，双凸时 $R_1,R_2$ 都取正）} $\dfrac1f=(n-1)\left(\dfrac1{R_1}+\dfrac1{R_2}\right)$。\ix{磨镜者公式}
几何光学符号法则写法：$\frac1f=(n-1)\left(\frac1{r_1}-\frac1{r_2}\right)$，双凸 $r_1>0,\ r_2<0$。

一般原则：\textbf{厚度 $\Delta(x,y)$ 的介质带来相位 $\varphi=k(n-1)\Delta(x,y)$}（相对空气，略去常数 $k\Delta_0$）。
\end{formula}

\begin{idea}[为什么“厚度 $\to$ 相位”是 $k(n-1)\Delta$]
光在介质里走 $\Delta$ 的路，相位涨 $kn\Delta$；同样这段路如果是空气，相位涨 $k\Delta$。和“全是空气”相比，多出 $k(n-1)\Delta$。
透镜中间厚、边缘薄，所以中间相位多、边缘相位少：把这个“中间多”的部分去掉常数，就剩边缘相对中间\textbf{少} $\frac{k r^2}{2f}$，即 $\exp(-ikr^2/2f)$。
\end{idea}

\begin{formula}{球面波的傍轴表示（判断会聚、发散、物像位置）}\ix{球面波!傍轴}
\begin{tabular}{@{}ll@{}}
从距离 $l$ 处点源发出的发散球面波，到达平面时 & $A\exp\!\left(+ik\dfrac{x^2+y^2}{2l}\right)$\\[6pt]
会聚到前方距离 $l'$ 处一点的球面波 & $A\exp\!\left(-ik\dfrac{x^2+y^2}{2l'}\right)$\\[6pt]
轴外点源 $(x_0,y_0)$、距离 $l$ & $A\exp\!\left[ik\dfrac{(x-x_0)^2+(y-y_0)^2}{2l}\right]$\\[6pt]
斜入射平面波（方向余弦 $\cos\alpha,\cos\beta$）& $A\exp[i2\pi(ux+vy)]$，$u=\dfrac{\cos\alpha}{\lambda}$
\end{tabular}

\smallskip
\textbf{读法}：二次相位系数 $\frac{k}{2R}$ 前为正 $\Rightarrow$ 发散，曲率中心在后方距离 $R$；为负 $\Rightarrow$ 会聚，在前方 $R$ 处成点。
$x$ 的一次项 $\exp(i2\pi u_0x)$ $\Rightarrow$ 波面倾斜（棱镜作用），会聚点横移 $\Delta x=\lambda f u_0$。
\end{formula}

\begin{recipe}{“某个相位分布等效于什么透镜？”（Alvarez、液晶透镜、超表面类题）}\ix{等效焦距}
\begin{enumerate}
\item 写出总相位 $\varphi(x,y)$（多个元件紧贴：相位相加；厚度型：$\varphi=k(n-1)\Delta$）。
\item 展开整理成 $\varphi=a(x^2+y^2)+bx+cy+\text{常数}$（若有三次项残留，就不是理想透镜）。
\item 对比 $-k\frac{x^2+y^2}{2F}$：$\boxed{F=-\dfrac{k}{2a}}$。$a<0$ 会聚，$a>0$ 发散。
\item 一次项 $bx$：$=2\pi u_0x$，$u_0=b/2\pi$，对应焦点横移 $F\lambda u_0=\dfrac{Fb}{k}$。
\item 屈光度 $\Phi=1/F$（单位 m$^{-1}$ = D）。
\end{enumerate}
\end{recipe}

\subsection*{单元：等效透镜题}\label{sec:equivlens}\ix{等效透镜}

\textbf{【题型】}“证明……等效于一个焦距为 $F$ 的透镜，并求 $F$”“$F$ 的正负与什么有关”“写出这个液晶/超表面的相位分布对应的焦距”。往年：2025-3(i)（Alvarez，见 §\ref{sec:alvarez}）。

\textbf{【完整推导（一般形式）】}设元件给出的厚度为 $\Delta(x,y)$（或直接给相位 $\varphi(x,y)$）。
\begin{enumerate}
\item 透过率：$t(x,y)=\exp[i\varphi(x,y)]$，$\varphi=k(n-1)\Delta(x,y)$。\quad 依据：厚度 $\to$ 相位（上面的直觉框）。
\item 平面波 $\ee^{ikz}$（在元件处为常数 1）照明，出射场 $=t(x,y)$。
\item 把 $\varphi$ 展开成多项式，按次数分组：常数项（略去，不影响强度和成像）、一次项（倾斜）、二次项（透镜）、三次及以上（像差）。
\item 若二次项为 $a(x^2+y^2)$，与理想透镜 $-\frac{k}{2F}(x^2+y^2)$ 逐项比较：$a=-\frac{k}{2F}\Rightarrow F=-\frac{k}{2a}$。
\item 判断：$F>0$ 出射为会聚球面波（会聚于后方 $F$ 处）；$F<0$ 为发散球面波，好像从前方 $|F|$ 处的点发出。
\end{enumerate}

\begin{exam}{答题模板：等效透镜}
解：元件的复透过率 $t(x,y)=\exp[ik(n-1)\Delta(x,y)]$。将 $\Delta$ 代入并整理：
\[
t(x,y)=\exp\!\Big[i\big(a(x^2+y^2)+bx+C\big)\Big],
\]
其中常数相位 $C$ 不影响结果，可略去。与薄透镜透过率 $\exp\!\big(-ik\frac{x^2+y^2}{2F}\big)$ 比较，得等效焦距
\[
F=-\frac{k}{2a}\quad(\text{代入 } a \text{ 的表达式}).
\]
$a<0$ 时 $F>0$，为会聚透镜；$a>0$ 时 $F<0$，为发散透镜。一次项 $bx$ 只使出射波面倾斜，焦点横移 $Fb/k$，不改变焦距。故该元件等效于焦距为 $F$ 的透镜。
\end{exam}

\textbf{【适配/变式】}
\begin{itemize}
\item \textbf{只有 $x^2$（柱透镜）}：$\varphi=ax^2$ $\Rightarrow$ 只在 $x$ 方向会聚，焦线平行于 $y$，$F_x=-k/2a$，$F_y=\infty$。$\varphi=a_xx^2+a_yy^2$ 且 $a_x\neq a_y$：像散透镜，两个焦距 $F_x=-k/2a_x$、$F_y=-k/2a_y$。
\item \textbf{两片紧贴}：相位相加 $\Rightarrow$ 屈光度相加 $\Phi=\Phi_1+\Phi_2$，即 $\frac1F=\frac1{f_1}+\frac1{f_2}$。
\item \textbf{偏心透镜} $\exp\!\big[-ik\frac{(x-x_0)^2+y^2}{2f}\big]$：展开 $=\exp\!\big(-ik\frac{x^2+y^2}{2f}\big)\exp\!\big(ik\frac{x_0x}{f}\big)\cdot$常数，即“居中透镜 + 棱镜”，焦点在 $(x_0,0,f)$。反过来，看到“二次项 + 一次项”就配方，配方的中心就是焦点横坐标。
\item \textbf{给的是相位而不是厚度}（液晶、SLM、超表面）：直接用 $\varphi$，不要再乘 $k(n-1)$。若写成 $\varphi=-\pi\frac{r^2}{\lambda F}$，就是 $-k\frac{r^2}{2F}$ 的另一种写法。
\item \textbf{给的是面型 $z(x,y)$ 而元件一面平}：厚度 $\Delta=\Delta_0-z$ 还是 $\Delta_0+z$ 要看凸起朝哪边；凸起处厚 $\Rightarrow$ 相位多。中心最厚 = 会聚。
\item \textbf{有三次及以上的项残留}：不是理想透镜，答题时指出“残留 $x^3$ 项（彗差型像差）/ $x^4$ 项（球差型）”。
\end{itemize}

\textbf{【数值例】}液晶透镜相位 $\varphi(r)=-\beta r^2$，$\beta=2\times10^{6}$ rad/m$^2$，$\lambda=0.55\,\mu$m。
$k=\frac{2\pi}{0.55\times10^{-6}}=1.142\times10^{7}$ m$^{-1}$，$a=-\beta$，$F=\frac{k}{2\beta}=\frac{1.142\times10^7}{4\times10^6}=2.86$ m，$\Phi=0.35$ D（会聚）。
同一元件换 $\lambda=0.65\,\mu$m：$F\propto k\propto1/\lambda$，$F=2.86\times\frac{0.55}{0.65}=2.42$ m（相位型元件有色散）。

\begin{warn}[等效透镜易错]
\begin{enumerate}
\item 二次项系数比较时别丢 $\frac12$：$-\frac{k}{2F}$，不是 $-\frac kF$。
\item 厚度型要乘 $(n-1)$，不是 $n$。
\item 屈光度单位：$F$ 用米，$\Phi=1/F$ 才是 D。
\end{enumerate}
\end{warn}

\needspace{8\baselineskip}\section{透镜的傅里叶变换性质}\label{sec:lensFT}

\begin{formula}{物在透镜前 $d_0$，平面波垂直照明，后焦面复振幅（课件核心结论）}\ix{透镜!傅里叶变换性质}\ix{傅里叶透镜}\ix{相位弯曲因子}
\[
\boxed{\tilde E(x,y)=\frac{\exp(ikf)}{i\lambda f}\exp\!\left[\frac{ik}{2f}\left(1-\frac{d_0}{f}\right)(x^2+y^2)\right]\F\{\tilde E_1(x_1,y_1)\}\Big|_{u=\frac{x}{\lambda f},\,v=\frac{y}{\lambda f}}}
\]
\begin{itemize}
\item $d_0=f$（物在前焦面，2f 系统）：二次相位消失，\textbf{精确傅里叶变换}。
\item $d_0=0$（紧贴透镜）：残留 $\exp\!\left(ik\frac{x^2+y^2}{2f}\right)$。
\item 任何 $d_0$：焦面\textbf{强度} $I=\dfrac1{\lambda^2f^2}\left|\F\{\tilde E_1\}\right|^2$ = 夫琅禾费衍射图样。
\item 坐标换算：焦面位置 $x=\lambda f\,u$。一个空间频率 $u_0$ 的分量聚焦在 $x=\lambda f u_0$ 处。
\item 推导：物$\to$透镜的自由空间传播在频域乘 $\exp[-i\pi\lambda d_0(u^2+v^2)]$；再乘紧贴情形的结论，代入 $u=x/\lambda f$ 即得。[F2-p14]
\end{itemize}
\end{formula}

\subsection*{完整推导：为什么后焦面是傅里叶变换 [F2-p08–15]}\label{sec:lensFTderiv}

\textbf{第 1 步（物紧贴透镜，$d_0=0$）}透镜后表面的场：
\[
\tilde E'(x_1,y_1)=\tilde E_1(x_1,y_1)\exp\!\left(-ik\frac{x_1^2+y_1^2}{2f}\right)\qquad\text{依据：透镜透过率 } t_L
\]
\textbf{第 2 步}后表面到后焦面（距离 $f$）做菲涅耳衍射：
\[
\tilde E(x,y)=\frac{\ee^{ikf}}{i\lambda f}\iint\tilde E'(x_1,y_1)\exp\!\left[ik\frac{(x-x_1)^2+(y-y_1)^2}{2f}\right]\dd x_1\dd y_1
\]
\textbf{第 3 步}展开平方 $(x-x_1)^2=x^2-2xx_1+x_1^2$，把只含 $x,y$ 的项提到积分外：
\[
\tilde E(x,y)=\frac{\ee^{ikf}}{i\lambda f}\ee^{ik\frac{x^2+y^2}{2f}}\iint\tilde E_1\,\underbrace{\ee^{-ik\frac{x_1^2+y_1^2}{2f}}\ee^{+ik\frac{x_1^2+y_1^2}{2f}}}_{=1\ (\text{透镜相位与菲涅耳相位抵消})}\exp\!\left(-ik\frac{xx_1+yy_1}{f}\right)\dd x_1\dd y_1
\]
\textbf{第 4 步}把 $\frac{k}{f}xx_1=2\pi\frac{x}{\lambda f}x_1$ 认作傅里叶核 $\ee^{-i2\pi ux_1}$，$u=\frac{x}{\lambda f}$：
\[
\tilde E(x,y)=\frac{\ee^{ikf}}{i\lambda f}\exp\!\left(ik\frac{x^2+y^2}{2f}\right)\F\{\tilde E_1\}\Big|_{u=\frac{x}{\lambda f},v=\frac{y}{\lambda f}}\qquad(d_0=0\text{ 的结论})
\]
\textbf{第 5 步（物在透镜前 $d_0$）}物到透镜前表面是自由传播 $d_0$。用第 3 章的角谱法：前表面的频谱 $=\F\{\tilde E_1\}\cdot\exp[-i\pi\lambda d_0(u^2+v^2)]$（略常数 $\ee^{ikd_0}$）。把它代替第 4 步中的 $\F\{\tilde E_1\}$。

\textbf{第 6 步}代入 $u=\frac{x}{\lambda f}$：$\pi\lambda d_0\frac{x^2+y^2}{\lambda^2f^2}=\frac{k}{2}\cdot\frac{d_0}{f^2}(x^2+y^2)$，与第 4 步的 $\frac{k}{2f}(x^2+y^2)$ 合并：
\[
\frac{k}{2f}(x^2+y^2)-\frac{k d_0}{2f^2}(x^2+y^2)=\frac{k}{2f}\left(1-\frac{d_0}{f}\right)(x^2+y^2)\quad\Rightarrow\quad\text{上面的方框公式}.
\]

\begin{formula}{物在透镜后方，距后焦面 $d$（会聚光照明）}\ix{透镜!物在透镜后}
\[
\tilde E(x,y)=\frac{\exp(ikd)}{i\lambda d}\cdot\frac{f}{d}\,\exp\!\left(ik\frac{x^2+y^2}{2d}\right)\F\{T(x_1,y_1)\}\Big|_{u=\frac{x}{\lambda d},\,v=\frac{y}{\lambda d}}
\]
（$f/d$ 是会聚球面波在屏上的振幅放大，课件略去了常数。）频谱尺度 $\lambda d$ 随屏的位置连续可调：屏越靠近焦面，频谱越小。[F2-p11]

推导一行：照到屏上的是会聚到焦点（距屏 $d$）的球面波 $\exp(-ik\frac{x_1^2+y_1^2}{2d})$；乘 $T$ 后传播 $d$ 做菲涅耳积分，菲涅耳的 $+ik\frac{x_1^2+y_1^2}{2d}$ 与它抵消，剩 $T$ 的傅里叶变换，频率 $u=\frac{x}{\lambda d}$。和第 1–4 步完全同构，只是 $f\to d$。
\end{formula}

\begin{formula}{球面波照明的一般情况}\ix{透镜!球面波照明}
点源距透镜 $l$，物在前焦面，在\textbf{点源的像面} $l'$（$\frac1l+\frac1{l'}=\frac1f$）观察：仍得物的精确傅里叶变换，频率坐标仍为 $u=x/\lambda f$。
关键一步用牛顿公式 $(l-f)(l'-f)=f^2$ 消掉二次相位。[F2-p16]

结论：\textbf{频谱面永远是点光源的像面}。平面波照明时点源在无穷远，像面就是后焦面。\ix{频谱面}
\end{formula}

\subsection*{单元：求焦面（频谱面）复振幅、强度、谱点位置}\label{sec:focalplane}

\textbf{【题型】}“衍射屏 $t(x_1,y_1)$ 置于透镜前 $d_0$ 处，单色平面波照明，求后焦面复振幅/光强分布”“频谱面上出现几个亮点，位置在哪”。

\textbf{【步骤】}
\begin{enumerate}
\item 写出紧贴衍射屏后的场 $\tilde E_1=$（照明光）$\times t(x_1,y_1)$。垂直平面波照明时照明光 = 1。
\item 查表求 $\F\{\tilde E_1\}(u,v)$（第 1 章变换对 + 性质：平移、缩放、卷积）。
\item 代入 $u=\frac{x}{\lambda f},v=\frac{y}{\lambda f}$，乘前面的常数和相位弯曲因子（$d_0=f$ 时没有）。
\item 求强度：取模平方，常数相位和二次相位全部消失：$I=\frac{1}{\lambda^2f^2}\big|\F\{\tilde E_1\}\big|^2_{u=x/\lambda f}$。
\item 谱点位置：频谱里的 $\delta(u-u_n)$ 在焦面上 $x_n=\lambda f u_n$ 处成亮点。
\end{enumerate}

\begin{exam}{答题模板：焦面复振幅（物在前焦面）}
解：衍射屏紧贴后的复振幅 $\tilde E_1(x_1,y_1)=t(x_1,y_1)$（单位振幅平面波垂直照明）。屏位于透镜前焦面，后焦面复振幅为屏函数的精确傅里叶变换：
\[
\tilde E(x,y)=\frac{\ee^{i2kf}}{i\lambda f}\,T\!\left(\frac{x}{\lambda f},\frac{y}{\lambda f}\right),\qquad T(u,v)=\F\{t(x_1,y_1)\}=\ \cdots(\text{查表写出}).
\]
光强 $I(x,y)=\dfrac{1}{\lambda^2f^2}\left|T\!\left(\dfrac{x}{\lambda f},\dfrac{y}{\lambda f}\right)\right|^2$。
（常数相位 $\ee^{i2kf}$ 是从屏到焦面的总光程 $2f$ 带来的，可略。）
\end{exam}

\textbf{【适配/变式】}
\begin{itemize}
\item \textbf{物不在前焦面（$d_0\neq f$）}：复振幅多乘 $\exp\!\big[\frac{ik}{2f}(1-\frac{d_0}{f})(x^2+y^2)\big]$；\textbf{强度不变}。问强度时直接说“二次相位因子取模后为 1”。
\item \textbf{斜入射平面波照明} $\ee^{i2\pi u_0x_1}$：$\tilde E_1=t\,\ee^{i2\pi u_0x_1}$，由调制定理频谱变为 $T(u-u_0,v)$，焦面图样\textbf{整体平移} $\lambda fu_0=f\cos\alpha$（$\approx f\sin\theta$）。
\item \textbf{屏平移} $t(x_1-a)$：频谱乘 $\ee^{-i2\pi ua}$，只是相位变化，\textbf{焦面强度不动}（夫琅禾费图样平移不变）。
\item \textbf{屏在透镜后}（距后焦面 $d$）：把 $\lambda f$ 全换成 $\lambda d$，多乘 $\frac fd$ 和二次相位 $\exp(ik\frac{x^2+y^2}{2d})$。
\item \textbf{点源（球面波）照明}：在点源的像面观察，频率坐标仍是 $x/\lambda f$。
\item \textbf{有限口径透镜}：只要口径足够大不挡住要用的频谱即可忽略；否则高于 $\frac{D/2}{\lambda f}$ 的频率被截掉（见 §\ref{sec:4fctf} 低通）。
\end{itemize}

\textbf{【数值例】}余弦振幅光栅 $t(x_1)=\frac12+\frac12\cos(2\pi x_1/d)$，$d=0.1$ mm，宽 $L=10$ mm（方孔），放在 $f=200$ mm 透镜的前焦面，$\lambda=632.8$ nm。
$T(u,v)=L^2\big[\frac12\sinc(Lu)+\frac14\sinc(L(u-\frac1d))+\frac14\sinc(L(u+\frac1d))\big]\sinc(Lv)$。
三个亮点：$x=0$ 和 $x=\pm\lambda f/d=\pm\frac{632.8\times10^{-6}\times200}{0.1}=\pm1.266$ mm；强度比 $1:\frac14:\frac14$（振幅比 $\frac12:\frac14:\frac14$）；每个亮点的宽度（到第一零点）$\lambda f/L=12.7\,\mu$m。

\needspace{8\baselineskip}\section{4f 系统}\label{sec:4f}

\begin{center}
\begin{tikzpicture}[scale=0.9,>=Stealth]
\draw[->] (-0.5,0)--(12.6,0) node[right]{$z$};
\foreach \x/\n in {0/输入面 $x_1$,4/频谱面 $\xi$,8/输出面 $x$} {\draw[thick] (\x*1.5,-1.3)--(\x*1.5,1.3); \node[below] at (\x*1.5,-1.3){\small \n};}
\foreach \x/\n in {2/$L_1(f_1)$,6/$L_2(f_2)$} {\draw[thick,fill=blue!10] (\x*1.5,0) ellipse (0.15 and 1.4); \node[above] at (\x*1.5,1.4){\small \n};}
\draw[<->] (0,-1.9)--(3,-1.9) node[midway,below]{$f_1$};
\draw[<->] (3,-1.9)--(6,-1.9) node[midway,below]{$f_1$};
\draw[<->] (6,-1.9)--(9,-1.9) node[midway,below]{$f_2$};
\draw[<->] (9,-1.9)--(12,-1.9) node[midway,below]{$f_2$};
\draw[red,thick] (0,0.8)--(3,0.8)--(6,0)--(9,-0.8*1.4)--(12,-0.8*1.4);
\node[red] at (1.5,1.05){\scriptsize 物点高 $h$};
\node[red,above] at (10.6,-1.05){\scriptsize 像高 $-\frac{f_2}{f_1}h$};
\end{tikzpicture}
\end{center}

\begin{formula}{4f 系统输出（无滤波）}\ix{4f 系统!输出}\ix{4f 系统!放大率}
频谱面：$\tilde E_f(\xi,\eta)=\dfrac{\ee^{i2kf_1}}{i\lambda f_1}\,\hat E\!\left(\dfrac{\xi}{\lambda f_1},\dfrac{\eta}{\lambda f_1}\right)$，$\hat E=\F\{\tilde E\}$。

输出面：
\[
\boxed{\tilde E_{\text{out}}(x,y)=-\,\ee^{i2k(f_1+f_2)}\,\frac{f_1}{f_2}\;\tilde E\!\left(-\frac{f_1}{f_2}x,\;-\frac{f_1}{f_2}y\right)},\qquad M=-\frac{f_2}{f_1}
\]
推导关键：$\displaystyle\int\exp\!\left[-i2\pi\xi\left(\frac{x_1}{\lambda f_1}+\frac{x}{\lambda f_2}\right)\right]\dd\xi=\delta\!\left(\frac{x_1}{\lambda f_1}+\frac{x}{\lambda f_2}\right)=\lambda f_1\,\delta\!\left(x_1+\frac{f_1}{f_2}x\right)$（二维为平方）。

振幅因子 $f_1/f_2=1/|M|$ 保证能量守恒：强度 $\times\frac1{M^2}$，面积 $\times M^2$。
\end{formula}

\subsection*{完整推导：4f 输出（两种方法，任选其一抄写）}\label{sec:4fderiv}

\textbf{方法一：二维直接积分 [F2-s2；F3-s80,82]}（保留振幅系数）
\begin{enumerate}
\item 第一个 2f（输入面在 $L_1$ 前焦面）：
\[
\tilde E_f(\xi,\eta)=\frac{\ee^{i2kf_1}}{i\lambda f_1}\iint\tilde E(x_1,y_1)\exp\!\left[-i2\pi\frac{\xi x_1+\eta y_1}{\lambda f_1}\right]\dd x_1\dd y_1.
\]
依据：透镜傅里叶变换性质，$d_0=f_1$。
\item 第二个 2f（频谱面在 $L_2$ 前焦面），仍是\textbf{正}变换：
\[
\tilde E_{\text{out}}(x,y)=\frac{\ee^{i2kf_2}}{i\lambda f_2}\iint\tilde E_f(\xi,\eta)\exp\!\left[-i2\pi\frac{x\xi+y\eta}{\lambda f_2}\right]\dd\xi\dd\eta.
\]
\item 代入并交换积分次序：
\[
\tilde E_{\text{out}}(x,y)=\frac{\ee^{i2k(f_1+f_2)}}{(i\lambda)^2f_1f_2}\iint\tilde E(x_1,y_1)\,K\dd x_1\dd y_1,
\]
其中 $K=\iint\exp\{-i2\pi[\xi(\frac{x_1}{\lambda f_1}+\frac{x}{\lambda f_2})+\eta(\frac{y_1}{\lambda f_1}+\frac{y}{\lambda f_2})]\}\dd\xi\dd\eta$。
\item 两个方向各用一次 $\int\ee^{-i2\pi\xi s}\dd\xi=\delta(s)$。设 $f_1,f_2>0$，由 $\delta(as)=\delta(s)/|a|$：
\[
K=(\lambda f_1)^2\delta\!\left(x_1+\frac{f_1}{f_2}x\right)\delta\!\left(y_1+\frac{f_1}{f_2}y\right).
\]
\item 筛选 $x_1=-\frac{f_1}{f_2}x,\ y_1=-\frac{f_1}{f_2}y$，前因子为 $\frac{(\lambda f_1)^2}{(i\lambda)^2f_1f_2}=-\frac{f_1}{f_2}$，故
\[
\tilde E_{\text{out}}(x,y)=-\ee^{i2k(f_1+f_2)}\frac{f_1}{f_2}\tilde E\!\left(-\frac{f_1}{f_2}x,-\frac{f_1}{f_2}y\right).
\]
\end{enumerate}

\textbf{方法二：占位符（尺度缩放）法 [F3-p81–82]}（多级、带滤波时更快）
\begin{enumerate}
\item 记 $\F\{\tilde E\}=\hat E(\blacksquare)$。第一次变换：频谱面场 $\propto\hat E\big(\frac{\xi}{\lambda f_1}\big)$。
\item 第二次变换把“$\hat E(\xi/\lambda f_1)$”当作 $\xi$ 的函数再做一次正变换。用 $\F\{\F\{g\}\}(\blacksquare)=g(-\blacksquare)$ 和缩放定理：
$\F_\xi\big\{\hat E(\tfrac{\xi}{\lambda f_1})\big\}(\nu)=\lambda f_1\,\tilde E(-\lambda f_1\nu)$。
\item 第二个透镜的频率是 $\nu=\frac{x}{\lambda f_2}$：$\tilde E_{\text{out}}\propto\tilde E\big(-\lambda f_1\frac{x}{\lambda f_2}\big)=\tilde E\big(-\frac{f_1}{f_2}x\big)$，$M=-\frac{f_2}{f_1}$。
\end{enumerate}
口诀：\textbf{每过一个 2f，把自变量“除以 $\lambda f$”填进占位符；两次正变换给一个负号。}

\begin{exam}{答题模板：4f 系统输出与放大率}
解：输入面位于 $L_1$ 前焦面，频谱面为 $L_1$ 后焦面，由透镜的傅里叶变换性质，频谱面复振幅
\[
\tilde E_f(\xi,\eta)=\frac{\ee^{i2kf_1}}{i\lambda f_1}\hat E\!\left(\frac{\xi}{\lambda f_1},\frac{\eta}{\lambda f_1}\right),\qquad\hat E=\F\{\tilde E\}.
\]
频谱面又位于 $L_2$ 前焦面，输出面为 $L_2$ 后焦面，再做一次傅里叶变换（频率 $\frac{x}{\lambda f_2},\frac{y}{\lambda f_2}$）。利用 $\F\{\F\{g\}\}(x)=g(-x)$ 及缩放定理，
\[
\tilde E_{\text{out}}(x,y)=-\ee^{i2k(f_1+f_2)}\frac{f_1}{f_2}\,\tilde E\!\left(-\frac{f_1}{f_2}x,-\frac{f_1}{f_2}y\right).
\]
即输出为输入的倒立像，横向放大率 $M=-\dfrac{f_2}{f_1}$；振幅系数 $\dfrac{f_1}{f_2}=\dfrac1{|M|}$ 保证能量守恒。
\end{exam}

\begin{formula}{4f 系统频谱面滤波（相干传递函数 CTF）}\ix{4f 系统!滤波}\ix{相干传递函数}\ix{CTF}
频谱面放复透过率 $T(\xi,\eta)$ 的滤波器。令物频率 $u=\xi/\lambda f_1$：
\[
\tilde E_{\text{out}}(x,y)=-\ee^{i2k(f_1+f_2)}\frac1{|M|}\Big[\tilde E\conv h\Big]\!\left(\frac xM,\frac yM\right),\qquad h=\IF\{T(\lambda f_1u,\lambda f_1v)\}
\]
\begin{itemize}
\item 以\textbf{物面}频率表示：$H(u,v)=T(\lambda f_1u,\ \lambda f_1v)$。
\item 以\textbf{像面}频率 $u'$ 表示：$u=Mu'$ 带符号，$H'(u',v')=T(\lambda f_1Mu',\lambda f_1Mv')=T(-\lambda f_2u',-\lambda f_2v')$。
\item 低通滤波器（圆孔半径 $a$）：物面截止频率 $u_c=\dfrac{a}{\lambda f_1}$，像面截止频率 $\dfrac{a}{\lambda f_2}$。
\item 若题目给的滤波器已写成 $T\!\left(\frac{\xi}{\lambda f_1},\frac{\eta}{\lambda f_1}\right)$，则物面 CTF 就是 $T(u,v)$。
\end{itemize}
\end{formula}

\subsection*{单元：频谱面滤波，求 CTF 和输出}\label{sec:4fctf}

\textbf{【题型】}“在频谱面设置复透射系数为 $T(\frac{\xi}{\lambda f_1},\frac{\eta}{\lambda f_1})$ 的滤波器，计算系统的相干传递函数”（2024-4(ii)）；“频谱面放半径 $a$ 的圆孔，求截止频率/输出”。

\textbf{【完整推导】}（占位符法；$\tilde E\to$ 输入，$\hat E=\F\{\tilde E\}$）
\begin{enumerate}
\item 频谱面滤波前：$\hat E\big(\frac{\xi}{\lambda f_1},\frac{\eta}{\lambda f_1}\big)$（略常数）。
\item 滤波后：$\hat E\big(\frac{\xi}{\lambda f_1},\cdot\big)\,T\big(\frac{\xi}{\lambda f_1},\cdot\big)=G\big(\frac{\xi}{\lambda f_1},\cdot\big)$，其中 $G(u,v)\equiv\hat E(u,v)T(u,v)$。\quad 依据：薄元件 = 乘透过率。
\item $L_2$ 正变换（同 4f 输出推导，只是把 $\hat E$ 换成 $G$）：$\tilde E_{\text{out}}(x,y)\propto g\big(-\frac{f_1}{f_2}x,-\frac{f_1}{f_2}y\big)=g\big(\frac xM,\frac yM\big)$，$g=\IF\{G\}=\tilde E\conv h$，$h=\IF\{T\}$。\quad 依据：卷积定理。
\item 求输出的频谱（以\textbf{物面}频率计，忽略倒像与缩放）：$\F\{g\}=\hat E(u,v)\,T(u,v)$。
\item 由 CTF 的定义（输出频谱/输入频谱）：$\boxed{H(u,v)=T(u,v)}$。若要以像面频率 $(u',v')$ 计：$u=Mu'$，$H'(u',v')=T(Mu',Mv')$。
\end{enumerate}

\begin{exam}{答题模板：4f 的相干传递函数}
解：频谱面复振幅为 $\hat E\big(\frac{\xi}{\lambda f_1},\frac{\eta}{\lambda f_1}\big)$（略去常数），经滤波器后变为
$\hat E\big(\frac{\xi}{\lambda f_1},\frac{\eta}{\lambda f_1}\big)T\big(\frac{\xi}{\lambda f_1},\frac{\eta}{\lambda f_1}\big)$。
$L_2$ 对其做傅里叶变换，输出为 $\tilde E_{\text{out}}(x,y)\propto\big[\tilde E\conv h\big]\big(\frac xM,\frac yM\big)$，$h=\F^{-1}\{T\}$，$M=-\frac{f_2}{f_1}$。
可见系统对物的作用是与 $h$ 卷积（相干点扩散函数），故以物面空间频率 $u=\frac{\xi}{\lambda f_1},v=\frac{\eta}{\lambda f_1}$ 表示的相干传递函数为
\[
H(u,v)=T(u,v);
\]
以像面频率 $(u',v')$ 表示，则 $u=Mu'$，$H'(u',v')=T(Mu',Mv')=T\!\left(-\frac{f_2}{f_1}u',-\frac{f_2}{f_1}v'\right)$。
\end{exam}

\textbf{【适配/变式】}
\begin{itemize}
\item \textbf{滤波器写成 $T(\xi,\eta)$（直接用位置）}：$H(u,v)=T(\lambda f_1u,\lambda f_1v)$。
\item \textbf{圆孔低通}（半径 $a$）：$T=\circf(\frac{r}{a})$，$H=\circf\big(\frac{\lambda f_1\rho}{a}\big)$，物面截止 $\rho_c=\frac{a}{\lambda f_1}$，像面截止 $\frac{a}{\lambda f_2}$。$\lambda=0.5\,\mu$m，$f_1=100$ mm，$a=1$ mm：$\rho_c=20$ mm$^{-1}$（能分辨周期 $\geq50\,\mu$m 的物结构）。
\item \textbf{狭缝}（沿 $y$ 方向的长缝，宽 $2a$）：只让 $|u|<\frac{a}{\lambda f_1}$ 通过，$v$ 不受限 $\Rightarrow$ 只模糊 $x$ 方向细节；竖条纹若频率超出就消失（阿贝-波特实验：狭缝方向与保留条纹方向\textbf{垂直} [F3-p78]）。
\item \textbf{两级级联、滤波器放在频谱面 2}（$L_3$ 后焦面，$T_2(\xi_2,\eta_2)$）：到达频谱面 2 前的场是中间像 $\tilde E_2(x)\propto\tilde E(-\frac{f_1}{f_2}x)$ 的频谱，$\hat E_2(\frac{\xi_2}{\lambda f_3})\propto\hat E\big(-\frac{f_2}{f_1}\cdot\frac{\xi_2}{\lambda f_3}\big)$。物频率 $u=-\frac{f_2\xi_2}{f_1\lambda f_3}$，即 $\xi_2=-\frac{\lambda f_1f_3}{f_2}u$：$H(u,v)=T_2\big(-\frac{\lambda f_1f_3}{f_2}u,-\frac{\lambda f_1f_3}{f_2}v\big)$。
\item \textbf{两个频谱面都放滤波器}：总 CTF 是两者（都换算到物频率后）的乘积。
\item \textbf{非相干照明}的 4f：不是本节内容，OTF = CTF 的归一化自相关（第 5 章）。
\end{itemize}

\begin{formula}{多级 4f 级联}\ix{4f 系统!级联}
$N$ 级：$M=\prod_j\left(-\dfrac{f_{2j}}{f_{2j-1}}\right)$。两级时 $M=\dfrac{f_2f_4}{f_1f_3}>0$，正立。
\[
\tilde E_{\text{out}}(x,y)=\ee^{i2k(f_1+f_2+f_3+f_4)}\frac{f_1f_3}{f_2f_4}\tilde E\!\left(\frac{f_1f_3}{f_2f_4}x,\frac{f_1f_3}{f_2f_4}y\right)
\]
推导：第一级输出（中间像）$\tilde E_2(x)=-\ee^{i2k(f_1+f_2)}\frac{f_1}{f_2}\tilde E(-\frac{f_1}{f_2}x)$；把 $\tilde E_2$ 当输入再套一次 4f 公式（$f_1\to f_3,f_2\to f_4$）：
$\tilde E_{\text{out}}(x)=-\ee^{i2k(f_3+f_4)}\frac{f_3}{f_4}\tilde E_2(-\frac{f_3}{f_4}x)$，两个负号相乘为正，自变量 $-\frac{f_1}{f_2}\cdot(-\frac{f_3}{f_4}x)=\frac{f_1f_3}{f_2f_4}x$（一维写法，二维振幅系数同样为 $\frac{f_1f_3}{f_2f_4}$）。完整答案见 2024-4（第 9 章）。
\end{formula}

\begin{formula}{频谱面常见操作对照表}\ix{空间滤波}\ix{相衬}\ix{4f 系统!线性相位}
\par\renewcommand{\arraystretch}{1.5}
\begin{tabular}{@{}P{4.6cm}P{5.6cm}P{5.2cm}@{}}
\toprule
频谱面放置 & 数学 & 输出效果\\ \midrule
圆孔光阑（低通）& 乘 $\circf(r/a)$ & 去掉细节，边缘模糊\\
挡住中心（高通）& 乘 $1-\circf(r/a)$ & 只剩边缘（边缘增强）\\
只留 $\pm1$ 级 & 只保留两个 $\delta$ & 余弦条纹，强度频率加倍\\
只留 0 和 +1 级 & & 余弦条纹，频率不变，对比度变化\\
挡掉 0 级（余弦光栅）& 去掉 $\delta(u)$ & 强度 $\propto\cos^2$：倍频条纹\\
线性相位 $\exp(i2\pi a\xi)$ & 频域乘线性相位 = 空域平移 & 像平移 $\Delta x=\lambda f_2\,a$（方向见下）\\
二次相位 $\exp(i\pi b\xi^2)$ & 等效频谱面加一薄透镜 & 改变像的离焦\\
相位板 $\ee^{i\pi/2}$ 盖住 0 级 & 泽尼克相衬 & 相位物体 $\to$ 强度可见\\
\bottomrule
\end{tabular}
\end{formula}

\subsection*{单元：输入光栅，频谱面挡掉某些级，求输出条纹}\label{sec:gratingfilter}\ix{空间滤波!光栅}\ix{余弦光栅!滤波}

\textbf{【题型】}“物为余弦（或矩形）振幅光栅，在频谱面放一狭缝/小屏，只让某几级通过，求像面光强、条纹周期和对比度，并画图”。

\textbf{【步骤】}（取 $f_1=f_2=f$，$M=-1$；$f_1\neq f_2$ 时把输出的 $x$ 换成 $x/M$）
\begin{enumerate}
\item 写出光栅透过率的傅里叶级数：$t(x)=\sum_nc_n\ee^{i2\pi nx/d}$。
\item 频谱 $\sum_nc_n\delta(u-\frac nd)$，第 $n$ 级亮点在频谱面 $\xi_n=\frac{n\lambda f}{d}$。
\item 判断哪些级被保留（比较 $\xi_n$ 与狭缝/光阑的尺寸）。
\item 输出复振幅 = 保留下来的那几项级数之和（倒像：$x\to-x$，对余弦对称光栅无影响）。
\item 求强度 $|\cdot|^2$，读出条纹频率和对比度 $V=\frac{I_{\max}-I_{\min}}{I_{\max}+I_{\min}}$。
\end{enumerate}

\begin{exam}{余弦光栅 $t(x)=\frac12+\frac m2\cos\frac{2\pi x}{d}$（$0<m\le1$），三种滤波}
\textbf{频谱}：$\frac12\delta(u)+\frac m4\delta(u-\frac1d)+\frac m4\delta(u+\frac1d)$，频谱面三点 $\xi=0,\pm\frac{\lambda f}{d}$。

\textbf{(a) 不滤波}：输出 $=t(-x)$，$I=\big(\frac12+\frac m2\cos\frac{2\pi x}d\big)^2$，周期 $d$。

\textbf{(b) 挡掉 0 级}（中心放小屏）：输出 $\frac m2\cos\frac{2\pi x}{d}$，
\[
I=\frac{m^2}{4}\cos^2\frac{2\pi x}{d}=\frac{m^2}{8}\Big(1+\cos\frac{4\pi x}{d}\Big)\quad\Rightarrow\quad\text{周期 }\frac d2\ \text{（倍频），对比度 }V=1.
\]
（只留 $\pm1$ 级结果相同。）

\textbf{(c) 只留 0 和 $+1$ 级}：输出 $\frac12+\frac m4\ee^{i2\pi x/d}$（倒像只改 $+1\leftrightarrow-1$，不影响强度），
\[
I=\Big|\frac12+\frac m4\ee^{i2\pi x/d}\Big|^2=\frac14+\frac{m^2}{16}+\frac m4\cos\frac{2\pi x}{d},\qquad V=\frac{m/4}{1/4+m^2/16}=\frac{4m}{4+m^2}.
\]
周期仍为 $d$；$m=1$ 时 $V=0.8$。（推导：$|a+b\ee^{i\theta}|^2=a^2+b^2+2ab\cos\theta$。）
\end{exam}

\begin{center}
\begin{tikzpicture}
\begin{axis}[width=12cm,height=4.2cm,xmin=0,xmax=2,ymin=0,ymax=1.05,samples=200,domain=0:2,
  xlabel={$x/d$},ylabel={$I$},xtick={0,0.5,1,1.5,2},ytick={0,0.25,0.5,1},
  legend style={font=\scriptsize,at={(1.02,1)},anchor=north west,draw=none},axis lines=left,tick label style={font=\scriptsize}]
\addplot[black,thick] {(0.5+0.5*cos(360*x))^2};
\addplot[red,thick,dashed] {0.25*(cos(360*x))^2};
\addplot[blue,thick,dotted] {0.3125+0.25*cos(360*x)};
\legend{{(a) 不滤波},{(b) 挡 0 级},{(c) 留 0 和 $+1$ 级}}
\end{axis}
\end{tikzpicture}

\small 余弦光栅 $m=1$ 三种滤波的输出强度。画图要点：横轴 $x$（标 $d/2,d$），(b) 的周期减半、最小值为 0；(c) 周期不变、最小值不为 0。
\end{center}

\textbf{【适配/变式】}
\begin{itemize}
\item \textbf{矩形（朗奇）光栅}，周期 $d$、透光宽 $a$：$c_n=\frac ad\sinc\frac{na}{d}$。$a=\frac d2$ 时 $c_0=\frac12$，$c_{\pm1}=\frac1\pi$，$c_{\pm2}=0$（缺级），$c_{\pm3}=-\frac1{3\pi}$。只留 $0,\pm1$ 级：输出 $\frac12+\frac2\pi\cos\frac{2\pi x}d$，方波被“抹圆”成余弦（注意 $\frac12-\frac2\pi<0$，强度在原暗区中间有一个小亮峰）。
\item \textbf{狭缝宽度 $w$ 的判断}：第 $n$ 级通过的条件 $|\xi_n|=\frac{|n|\lambda f}{d}<\frac w2$。
\item \textbf{二维正交网格} + 水平狭缝：只留 $v=0$ 那一行谱点 $\Rightarrow$ 输出只剩 $x$ 方向变化，即\textbf{竖直}条纹；竖直狭缝 $\Rightarrow$ 水平条纹 [F3-p78]。
\item \textbf{$f_1\neq f_2$}：输出条纹周期乘 $|M|=\frac{f_2}{f_1}$。
\item \textbf{相位光栅} $t=\ee^{i\phi_0\cos(2\pi x/d)}$：不滤波时强度均匀（看不见）；挡掉 0 级或用相衬才能看见（§\ref{sec:processing}）。
\end{itemize}

\begin{recipe}{频谱面放线性相位后输出平移多少}
设频谱面相位 $\varphi(\xi)=2\pi a\xi$（例如 $kQ\xi$，则 $a=Q/\lambda$）。第二个透镜做正变换，求值于 $\nu=x/\lambda f_2$：
\[
\F\{G(\xi)\ee^{i2\pi a\xi}\}(\nu)=\hat G(\nu-a)\quad\Rightarrow\quad \text{输出在 } x=\lambda f_2a \text{ 处的值等于原来 } x=0 \text{ 处的值}
\]
即\textbf{输出整体沿 $+x$ 平移 $\Delta x=\lambda f_2a$}。$\varphi=kQ\xi$ 时 $\Delta x=Qf_2$（几何意义：频谱面加棱镜，偏转角 $Q$）。
无平移时 4f 输出为 $t_1(-x)$，有平移时为 $t_1(-(x-\Delta x))=t_1(\Delta x-x)$。
\end{recipe}

\subsection*{单元：频谱面加线性相位（棱镜、SLM 倾斜相位）}\label{sec:linphase}\ix{线性相位!频谱面}

\textbf{【题型】}“频谱面放置 SLM，相位 $\varphi(\xi,\eta)=kQ\xi$，求输出复振幅”（2025-3(ii) 的第一步）；“频谱面放一个小角度棱镜，像移动多少”。

\textbf{【完整推导】}（$f_1=f_2=f$；一般情形把第二个透镜的 $f$ 换成 $f_2$）
\begin{enumerate}
\item 频谱面加相位前：$\hat t\big(\frac{\xi}{\lambda f}\big)$（输入 $t(x)$，$\hat t=\F\{t\}$）。
\item 乘相位：$\hat t\big(\frac{\xi}{\lambda f}\big)\exp(ikQ\xi)=\hat t\big(\frac{\xi}{\lambda f}\big)\exp\!\big(i2\pi\frac{Q}{\lambda}\xi\big)$。
\item $L_2$ 正变换，频率 $\nu=\frac{x'}{\lambda f}$：
$\displaystyle\int\hat t\Big(\frac{\xi}{\lambda f}\Big)\exp\!\Big[-i2\pi\xi\Big(\frac{x'}{\lambda f}-\frac Q\lambda\Big)\Big]\dd\xi=\int\hat t\Big(\frac{\xi}{\lambda f}\Big)\exp\!\Big[-i2\pi\xi\,\frac{x'-Qf}{\lambda f}\Big]\dd\xi$。
\item 与无相位时的结果（$t(-x')$）比较，只是 $x'\to x'-Qf$：输出 $=t\big(-(x'-Qf)\big)=t(Qf-x')$。\quad 依据：调制定理（频域乘线性相位 $\Leftrightarrow$ 空域平移）。
\item 结论：像沿 $+x'$ 方向平移 $Qf$（$f_1\neq f_2$ 时为 $Qf_2$），像本身仍是倒像。
\end{enumerate}

\textbf{【适配/变式】}
\begin{itemize}
\item $\varphi=kQ_x\xi+kQ_y\eta$：二维平移 $(Q_xf_2,\ Q_yf_2)$。
\item $\varphi=-kQ\xi$：反向平移，输出 $t(-Qf-x')$。
\item 写成 $\varphi=2\pi a\xi$：$\Delta x=\lambda f_2a$；写成光栅频率 $\varphi=2\pi\xi/\Lambda$：$\Delta x=\lambda f_2/\Lambda$。
\item \textbf{对偶情形}：线性相位放在\textbf{输入面}（斜入射照明 $\ee^{i2\pi u_0x}$）$\Rightarrow$ 频谱在频谱面平移 $\lambda f_1u_0$，输出像不动、只多一个线性相位。
\item 数值：$\lambda$ 无关！$Q=0.01$ rad，$f=100$ mm：$\Delta x=1$ mm。
\end{itemize}

\needspace{8\baselineskip}\section{物不在前焦面时的 4f 系统（纵向偏离）}\label{sec:4foffset}

\begin{formula}{偏离补偿：$x_2=M^2x_1$}\ix{4f 系统!离焦补偿}\ix{纵向放大率}
物面在前焦面\textbf{外侧} $x_1$（距 $L_1$ 为 $f_1+x_1$），在后焦面\textbf{前方} $x_2$（距 $L_2$ 为 $f_2-x_2$）接收。

频域看：物多走了 $x_1$，频谱面多出传播因子 $\exp[-i\pi\lambda x_1(u^2+v^2)]$。
换成输出面频率 $u'=-\dfrac{f_1}{f_2}u$ 后为 $\exp\!\left[-i\pi\lambda x_1\dfrac{f_2^2}{f_1^2}(u'^2+v'^2)\right]$。
从后焦面往回走 $x_2$，乘 $\exp[+i\pi\lambda x_2(u'^2+v'^2)]$。两者抵消的条件：
\[
\boxed{x_2=\left(\frac{f_2}{f_1}\right)^2x_1=M^2x_1}
\]
这就是几何光学的\textbf{纵向放大率} $\alpha=M^2$（视光学 [VO-p35]）。
\end{formula}

\subsection*{单元：证明 $x_2=M^2x_1$（2024-5）}

\textbf{【题型】}“输入复振幅偏离 $L_1$ 前焦面 $x_1$，证明在 $L_2$ 后焦面前 $x_2$ 处接收仍能实现 4f 的作用，$x_2=M^2x_1$”。逐步完整答案见第 9 章 2024-5（第 \pageref{ex:2024-5} 页），这里给思路骨架与变式。

\textbf{【思路骨架】}整个系统都在\textbf{频域}算，每一段都是“乘一个因子”：
\begin{enumerate}
\item 输入面 $\to$ 前焦面：自由传播 $x_1$，频谱乘 $H_{x_1}(u)=\exp[-i\pi\lambda x_1u^2]$（一维写；二维 $u^2\to u^2+v^2$）。
\item 前焦面 $\to$ 后焦面：标准 4f，输出 $\propto g(\frac xM)$，其频谱以输出频率 $u'$ 计为 $\hat g(Mu')$（缩放定理，振幅常数略）。所以后焦面场的频谱 $=\hat E(Mu')\exp[-i\pi\lambda x_1M^2u'^2]$。
\item 后焦面 $\to$ 接收面（在其前方 $x_2$）：等于“反向传播 $x_2$”，乘 $H_{-x_2}(u')=\exp[+i\pi\lambda x_2u'^2]$。
\item 接收面频谱 $=\hat E(Mu')\exp[-i\pi\lambda(M^2x_1-x_2)u'^2]$。取 $x_2=M^2x_1$ 时二次相位消失，接收面场 $\propto\tilde E(\frac xM)$，与标准 4f 输出相同。证毕。
\end{enumerate}

\textbf{【几何光学验算】}（可写在答案最后加分）用牛顿公式 $xx'=f^2$（物、像到各自焦点的距离，物在焦点外为正）：
$L_1$：物距前焦点 $x_1$ $\Rightarrow$ 中间像在 $L_1$ 后焦点外 $\frac{f_1^2}{x_1}$ 处；
该点相对 $L_2$ 的前焦点（与 $L_1$ 后焦点重合）在\textbf{靠近 $L_2$} 一侧 $\frac{f_1^2}{x_1}$，即物距 $=-\frac{f_1^2}{x_1}$；
$L_2$：像距后焦点 $x'=\frac{f_2^2}{-f_1^2/x_1}=-\frac{f_2^2}{f_1^2}x_1$，负号表示在后焦点\textbf{之前}，距离 $M^2x_1$。与波动光学结果一致。

\textbf{【适配/变式】}
\begin{itemize}
\item 物在前焦面\textbf{内侧}（$x_1<0$）：$x_2=M^2x_1<0$，接收面在后焦面\textbf{之后} $M^2|x_1|$。
\item $f_1=f_2$：$x_2=x_1$，物往外移多少，接收面就往里移多少。
\item 只移动接收面、不移物（$x_1=0,x_2\neq0$）：残留 $\exp[i\pi\lambda x_2u'^2]$，即离焦，像模糊（强度不再是 $|E|^2$ 的缩放）。
\item 问横向放大率是否改变：不变，仍为 $M=-\frac{f_2}{f_1}$（频率换算只用到 $f_1,f_2$）。
\end{itemize}

\needspace{8\baselineskip}\section{Alvarez / Lohmann 变焦透镜}\label{sec:alvarez}

\begin{formula}{三次面型平移 $\Rightarrow$ 二次相位}\ix{Alvarez 透镜}
面型 $T(x,y)=A\left(\frac13x^3+xy^2\right)+B$，$L_1$ 厚度 $T$，$L_2$ 厚度 $-T$（互补）。

$L_1$ 平移 $-d$、$L_2$ 平移 $+d$（沿 $x$ 反向，各移 $d$）：总厚度
\[
\Delta=T(x+d,y)-T(x-d,y)=2Ad(x^2+y^2)+\tfrac23Ad^3\quad(\text{三次多项式，结果\textbf{精确}})
\]
（“导数的原始定义”：$\frac{T(x+d)-T(x-d)}{2d}\to\frac{\partial T}{\partial x}=A(x^2+y^2)$。）

相位 $\varphi=k(n-1)\Delta=2k(n-1)Ad\,(x^2+y^2)+\text{常数}$。对比 $-k\frac{x^2+y^2}{2F}$：
\[
\boxed{F=-\frac{1}{4(n-1)Ad}}
\]
$Ad>0$ 时为发散透镜；平移方向反过来（$d\to-d$）就变成会聚透镜，$F=+\frac1{4(n-1)Ad}$。$|F|\propto1/d$，平移越多，焦距越短。
\end{formula}

\subsection*{单元：2025-3(i) Alvarez 平移变焦，逐步解}\label{sec:alvarezi}

\textbf{【原题要点】}$L_1$ 面型 $T(x,y)=A(\frac13x^3+xy^2)+B$，$L_2$ 面型 $-T$，薄透镜、间距为零，可直接用厚度算相位；两片沿 $x$ \textbf{相向}平移同样距离 $d$（图：$L_1$ 向下即 $-x$，$L_2$ 向上即 $+x$）。证明等效于焦距 $F$ 的透镜，$F$ 的正负与平移方向有关，求 $F$。提示：导数的原始定义。

\textbf{【完整推导】}
\begin{enumerate}
\item 平移后的厚度：$L_1$ 沿 $-x$ 移 $d$，原来在 $x$ 处的厚度现在在 $x-d$ 处，所以新厚度 $T(x+d,y)$；$L_2$ 沿 $+x$ 移 $d$，新厚度 $-T(x-d,y)$。\quad 依据：函数 $f(x)$ 右移 $a$ 得 $f(x-a)$。
\item 两片紧贴，厚度相加：$\Delta(x,y)=T(x+d,y)-T(x-d,y)$（$B$ 抵消）。
\item 展开立方：$(x+d)^3-(x-d)^3=6x^2d+2d^3$；$(x+d)y^2-(x-d)y^2=2dy^2$。所以
\[
\Delta=A\Big[\tfrac13(6x^2d+2d^3)+2dy^2\Big]=2Ad(x^2+y^2)+\tfrac23Ad^3.
\]
（提示的用法：$\Delta=2d\cdot\frac{T(x+d)-T(x-d)}{2d}\approx2d\,\frac{\partial T}{\partial x}=2Ad(x^2+y^2)$，$d$ 小时成立；直接展开发现对三次面型\textbf{精确}成立，只多一个常数。）
\item 相位：$\varphi=k(n-1)\Delta=2k(n-1)Ad(x^2+y^2)+\frac23k(n-1)Ad^3$，常数项略去。
\item 比较 $-\frac{k}{2F}(x^2+y^2)$：$-\frac{k}{2F}=2k(n-1)Ad\Rightarrow F=-\dfrac{1}{4(n-1)Ad}$。
\item 讨论符号（设 $A>0$，$n>1$）：按图示方向 $d>0$，$F<0$，发散；两片反向平移（$d<0$）则 $F>0$，会聚。屈光度 $\Phi=-4(n-1)Ad$ 与 $d$ 成正比，连续可调。
\end{enumerate}

\begin{exam}{答题模板：2025-3(i)}
解：$L_1$ 沿 $-x$ 平移 $d$ 后厚度为 $T(x+d,y)$，$L_2$ 沿 $+x$ 平移 $d$ 后厚度为 $-T(x-d,y)$。两薄透镜紧贴，总厚度
\[
\Delta(x,y)=T(x+d,y)-T(x-d,y)=A\Big[\tfrac13\big((x+d)^3-(x-d)^3\big)+2dy^2\Big]=2Ad(x^2+y^2)+\tfrac23Ad^3 .
\]
（由导数定义，$\Delta=2d\cdot\frac{T(x+d,y)-T(x-d,y)}{2d}\approx2d\,\partial_xT=2Ad(x^2+y^2)$，对三次面型精确成立。）
附加相位 $\varphi=k(n-1)\Delta=2k(n-1)Ad\,(x^2+y^2)+$常数。与薄透镜相位 $-\frac{k}{2F}(x^2+y^2)$ 比较，
\[
F=-\frac{1}{4(n-1)Ad}.
\]
当 $Ad>0$（按图示方向平移）时 $F<0$，为发散透镜；平移方向相反（$d\to-d$）时 $F>0$，为会聚透镜。$d=0$ 时 $F\to\infty$，无光焦度。
\end{exam}

\textbf{【适配/变式】}
\begin{itemize}
\item \textbf{只移一片}，相对位移 $s$（$L_1$ 不动，$L_2$ 沿 $+x$ 移 $s$）：$\Delta=T(x,y)-T(x-s,y)=A\big[s(x^2+y^2)-s^2x+\frac{s^3}3\big]$，$F=-\frac{1}{2(n-1)As}$，另有一次项 $-As^2x$（焦点横移 $\frac s2$）。\textbf{只看相对位移}：对称各移 $d$ 等价于 $s=2d$，且对称平移没有一次项。
\item \textbf{沿 $y$ 平移}：$T(x,y+d)-T(x,y-d)=4Adxy$，是马鞍面（像散），\textbf{不是}透镜。所以平移方向必须沿 $x$。
\item \textbf{面型只有 $\frac A3x^3$}：$\Delta=2Adx^2+$常数，柱透镜，只在 $x$ 方向有焦距 $F_x=-\frac1{4(n-1)Ad}$。
\item \textbf{面型系数写法不同}（如 $T=a(x^3+3xy^2)$，即 $A=3a$）：先化成 $A(\frac13x^3+xy^2)$ 再代公式。
\end{itemize}

\textbf{【数值例】}$n=1.5$，$A=1000$ m$^{-2}$（在 $x=5$ mm 处面型高 $\frac A3x^3\approx42\,\mu$m），$d=1$ mm：
$F=-\frac{1}{4\times0.5\times1000\times10^{-3}}=-0.5$ m，$\Phi=-2$ D；$d=2$ mm 时 $\Phi=-4$ D。即每平移 1 mm 改变 $-2$ D。

\begin{formula}{4f 实现 Alvarez（Split-Lohmann）：输入面放 $L_1$，频谱面放线性相位，输出面放 $L_2$}\ix{Alvarez 透镜!4f 实现}\ix{Split-Lohmann}
\begin{enumerate}
\item 输入：$t_1(x,y)=\exp[ik(n-1)T(x,y)]$。
\item 频谱面乘 $\exp(ikQ\xi)$：由上面的套路，4f 输出（$f_1=f_2=f$，$M=-1$）为 $t_1(Qf-x',\,-y')$。
\item 再乘 $L_2$：$t_2(x',y')$。总输出 $\tilde E_{\text{out}}=t_1(Qf-x',-y')\cdot t_2(x',y')$（乘常数 $-\ee^{i4kf}$）。
\item \textbf{关键}：倒像 $(x,y)\to(-x,-y)$ 把奇函数 $T$ 变成 $-T$：$T(-x,-y)=-T(x,y)+2B$。所以
\[
t_1(Qf-x',-y')=\exp\{-ik(n-1)[T(x'-Qf,y')-2B]\}
\]
\item 若 $L_2$ 厚度为 $+T$（与 $L_1$ 同型，倒像后正好互补）：
\[
\tilde E_{\text{out}}\propto\exp\{ik(n-1)[T(x',y')-T(x'-Qf,y')]\}
=\exp\!\Big\{ik(n-1)A\big[Qf(x'^2+y'^2)-Q^2f^2x'\big]\Big\}
\]
二次项给出焦距 $\boxed{F'=-\dfrac{1}{2(n-1)AQf}}$；一次项 $-Q^2f^2x'$ 是倾斜，只让焦点横移。与平移版对照：相对平移量 $2d\leftrightarrow Qf$。
\item 若 $L_2$ 按题面取 $-T$，则三次项不抵消：
$\tilde E\propto\exp\{-ik(n-1)A[\tfrac23x'^3+2x'y'^2-Qf(x'^2+y'^2)+Q^2f^2x']\}$，残留三次相位，不是理想透镜。
答题时写出这一步并说明：要实现变焦，第二片在 4f 后应取与第一片同号的面型（因为 4f 倒像）。
\end{enumerate}
\end{formula}

\subsection*{单元：2025-3(ii) 4f 版 Alvarez，逐步解}\label{sec:alvarezii}

\textbf{【原题要点】}4f 系统（两透镜焦距都是 $f$），$L_1$ 放输入面，$L_2$ 放输出面，频谱面放 SLM，相位 $\varphi(\xi,\eta)=kQ\xi$；三个面的坐标 $x,\xi,x'$ 同向（都朝上）。求 $L_2$ 右表面复振幅（系统输出）和焦距 $F'$。

\textbf{【完整推导】}
\begin{enumerate}
\item 输入面：单位平面波经 $L_1$，$t_1(x,y)=\exp[ik(n-1)T(x,y)]$。
\item 频谱面：$\hat t_1\big(\frac{\xi}{\lambda f},\frac{\eta}{\lambda f}\big)$，乘 SLM 相位 $\exp(ikQ\xi)$。
\item $L_2$ 前的场（4f 输出 + 线性相位平移，见 §\ref{sec:linphase}）：$-\ee^{i4kf}\,t_1(Qf-x',-y')$。
\item 用奇函数性质 $T(-X,-Y)=-T(X,Y)+2B$，取 $X=x'-Qf$，$Y=y'$：$t_1(Qf-x',-y')=\exp\{ik(n-1)[-T(x'-Qf,y')+2B]\}$。
\item 乘 $L_2$ 的透过率 $t_2=\exp[ik(n-1)T_2(x',y')]$，输出
$\tilde E_{\text{out}}=-\ee^{i4kf}\ee^{2ik(n-1)B}\exp\{ik(n-1)[T_2(x',y')-T(x'-Qf,y')]\}$。
\item \textbf{情况 A}：$T_2=+T$。$T(x',y')-T(x'-s,y')$（$s=Qf$）$=A\big[\frac13(x'^3-(x'-s)^3)+sy'^2\big]=A\big[s(x'^2+y'^2)-s^2x'+\frac{s^3}{3}\big]$。
\[
\tilde E_{\text{out}}\propto\exp\!\Big\{ik(n-1)A\big[Qf(x'^2+y'^2)-Q^2f^2x'\big]\Big\}
\]
比较 $-\frac{k}{2F'}$：$\boxed{F'=-\dfrac{1}{2(n-1)AQf}}$。一次项使焦点横移 $\frac{Qf}{2}$（套路第 4 步：$F'\cdot\frac{b}{k}$，$b=-k(n-1)AQ^2f^2$）。
\item \textbf{情况 B}：$T_2=-T$（题面字面）。$-T(x')-T(x'-s)$ 中三次项 $-\frac{2A}{3}x'^3-2Ax'y'^2$ 留下，不是理想透镜。
\end{enumerate}

\begin{exam}{答题模板：2025-3(ii)}
解：输入面复振幅 $t_1(x,y)=\exp[ik(n-1)T(x,y)]$。频谱面复振幅为 $\hat t_1\big(\frac{\xi}{\lambda f},\frac{\eta}{\lambda f}\big)$（略常数），乘 SLM 相位 $\exp(ikQ\xi)=\exp(i2\pi\frac{Q}{\lambda}\xi)$ 后，经 $L_2$ 前的透镜再做傅里叶变换（频率 $\frac{x'}{\lambda f}$），由平移定理得
\[
\tilde E(x',y')=-\ee^{i4kf}\,t_1(Qf-x',-y'),
\]
即倒像并沿 $x'$ 平移 $Qf$。由于 $T(-X,-Y)=-T(X,Y)+2B$，
$t_1(Qf-x',-y')=\exp\{-ik(n-1)[T(x'-Qf,y')-2B]\}$。
倒像使 $L_1$ 的作用等效于面型 $-T$，故 $L_2$ 取面型 $+T$ 时，$L_2$ 右表面复振幅为
\[
\tilde E_{\text{out}}(x',y')=C\exp\Big\{ik(n-1)A\big[Qf(x'^2+y'^2)-Q^2f^2x'\big]\Big\},
\]
$C$ 为常数。二次项对应焦距 $F'=-\dfrac{1}{2(n-1)AQf}$，一次项只使焦点横移，不影响焦距。改变 $Q$ 即可变焦（$Qf$ 相当于平移版中的相对位移 $2d$）。
若 $L_2$ 面型取 $-T$，输出相位为 $-k(n-1)[T(x'-Qf,y')+T(x',y')]$，三次项不能抵消，系统不是理想透镜。
\end{exam}

\textbf{【数值例】}$n=1.5$，$A=1000$ m$^{-2}$，$f=0.1$ m，要 $F'=-0.5$ m：$Q=\frac{1}{2\times0.5\times1000\times0.1\times0.5}=0.02$ rad，$Qf=2$ mm，与平移版 $d=1$ mm（$2d=2$ mm）给出同样焦距。

\begin{warn}[4f 题最常丢分的地方]
\begin{enumerate}
\item 忘了 4f 输出是\textbf{倒像} $\tilde E(-x/|M|)$，以及振幅系数 $1/|M|$。
\item 滤波器的坐标：频谱面位置 $\xi$ 和物频率 $u$ 之间差 $\lambda f_1$；放大后像面频率再差 $M$。
\item 第二个透镜做的是\textbf{正}变换（不是逆变换），所以出现 $f(-x)$。
\item 常数相位 $\ee^{ik\cdot\text{总光程}}$ 和 $1/(i\lambda f)$ 可写可略，但要说明“略去常数因子”。
\item 平移方向：频谱面相位 $+kQ\xi$ 使输出沿 $+x'$ 移 $Qf$；先写 $t(-x')$ 再把 $x'$ 换成 $x'-Qf$，不要先平移再倒像。
\item 奇函数倒像变号（Alvarez）、偶函数倒像不变：倒像前先看函数的奇偶性。
\end{enumerate}
\end{warn}

\needspace{8\baselineskip}\section{相干光学处理（课件 III）}\label{sec:processing}
\begin{formula}{4f 的 CTF 就是频谱面滤波器 [F3-p84]}\ix{4f 系统!CTF 证明}
$f(x,y)\to h(\xi,\eta)=F\!\left(\frac{\xi}{\lambda f},\frac{\eta}{\lambda f}\right)\to h'=F\cdot T\to g(x',y')=H'\!\left(\frac{x'}{\lambda f},\frac{y'}{\lambda f}\right)$，
输出频谱 $G(u,v)=h'(-\lambda fu,-\lambda fv)=F(-u,-v)T(-u,-v)$。忽略 $-1$ 倒像：$G=F\cdot T$，即 \textbf{CTF $=T(u,v)$}。
低通 $\to$ 模糊；高通（挡直流）$\to$ 边缘；带通（圆环）$\to$ 一定尺度的轮廓 [F3-p85]。

（课件这里 $F$ 是物 $f$ 的频谱，$H'$ 是 $h'$ 的频谱，$\blacksquare$ 占位符法：第三次变换求输出频谱时，$g(x')=H'(\frac{x'}{\lambda f})$ 再变换得 $h'(-\lambda f\blacksquare)$。）
\end{formula}

\begin{idea}[高通为什么得到边缘]
平坦区域只有低频（接近直流），边缘处亮度突变，含大量高频。挡住中心（低频）后，平坦区域变暗，只有边缘还有光。
低通相反：把高频（细节、锐边）挡掉，图像变模糊。带通（圆环）只留某一尺度的结构。
\end{idea}

\begin{formula}{光学加减法：频谱面放余弦光栅 [F3-p86–88]}\ix{光学加减法}\ix{余弦光栅!频谱面}
输入 $E(x)=E_1(x-b)+E_2(x+b)$（两个图形相距 $2b$），频谱面 $T=\frac12+\frac12\cos(2\pi bu)$，$u=\frac{\xi}{\lambda f}$。
\[
\big[\mathscr E_1\ee^{-i2\pi ub}+\mathscr E_2\ee^{i2\pi ub}\big]\Big[\tfrac12+\tfrac14\ee^{i2\pi bu}+\tfrac14\ee^{-i2\pi bu}\Big]
\]
展开、逆变换（乘 4、忽略倒像）：
\[
g(x')=\boxed{E_1(x')+E_2(x')}+2E_1(x'-b)+E_1(x'-2b)+2E_2(x'+b)+E_2(x'+2b)
\]
光栅把每个图形复制成 0、$\pm1$ 级三份，间距 $b$；$E_1$ 的 $+1$ 级与 $E_2$ 的 $-1$ 级落在中心相加。（课件此处 $E_1(x'-2b)$ 的系数写成 2，按推导应为 1。）

\textbf{变减法}：光栅取 $\frac12+\frac12\cos(2\pi bu+\phi_0)$，中心项变为 $\frac14\ee^{i\phi_0}[E_1+\ee^{-2i\phi_0}E_2]$。
要 $\ee^{-2i\phi_0}=-1$ 需 $\phi_0=\frac\pi2$，即光栅横向平移\textbf{四分之一周期}（课件说“平移 $\pi/4$”，指的就是这个；答题写清是相位 $\pi/2$ 即 1/4 周期）。
\end{formula}

\subsection*{单元：光学加减法逐步推导与答题模板}\label{sec:addsub}

\textbf{【完整推导】}（记 $\mathscr E_j=\F\{E_j\}$，$u=\frac{\xi}{\lambda f}$）
\begin{enumerate}
\item 输入频谱：$\F\{E_1(x-b)\}=\mathscr E_1(u)\ee^{-i2\pi ub}$，$\F\{E_2(x+b)\}=\mathscr E_2(u)\ee^{+i2\pi ub}$。\quad 依据：平移定理。
\item 光栅写成指数：$\frac12+\frac12\cos(2\pi bu)=\frac12+\frac14\ee^{i2\pi bu}+\frac14\ee^{-i2\pi bu}$。\quad 依据：欧拉公式。
\item 相乘得六项（每项的指数相加）：
\[
\tfrac12\mathscr E_1\ee^{-i2\pi bu}+\tfrac14\mathscr E_1+\tfrac14\mathscr E_1\ee^{-i4\pi bu}
+\tfrac12\mathscr E_2\ee^{i2\pi bu}+\tfrac14\mathscr E_2\ee^{i4\pi bu}+\tfrac14\mathscr E_2 .
\]
\item 逐项变回空域（$\mathscr E\ee^{-i2\pi uc}\to E(x'-c)$，忽略倒像和常数）：
$\frac12E_1(x'-b)+\frac14E_1(x')+\frac14E_1(x'-2b)+\frac12E_2(x'+b)+\frac14E_2(x'+2b)+\frac14E_2(x')$。
\item 在 $x'=0$ 附近（原两图的中点）：$\frac14[E_1(x')+E_2(x')]$——两图重合相加。其余四项在 $\pm b,\pm2b$ 处，不重叠（要求图形尺寸 $<b$）。
\end{enumerate}

\begin{exam}{答题模板：用 4f 系统实现两图相减}
解：输入 $E(x)=E_1(x-b)+E_2(x+b)$，频谱为 $\mathscr E_1(u)\ee^{-i2\pi bu}+\mathscr E_2(u)\ee^{i2\pi bu}$，$u=\frac{\xi}{\lambda f}$。
频谱面放余弦振幅光栅 $T=\frac12+\frac12\cos(2\pi bu+\phi_0)=\frac12+\frac14\ee^{i\phi_0}\ee^{i2\pi bu}+\frac14\ee^{-i\phi_0}\ee^{-i2\pi bu}$，相乘后不含指数因子（落在像面中心）的项为
\[
\tfrac14\ee^{i\phi_0}\mathscr E_1(u)+\tfrac14\ee^{-i\phi_0}\mathscr E_2(u),
\]
对应像面中心复振幅 $\frac14\ee^{i\phi_0}\big[E_1(x')+\ee^{-i2\phi_0}E_2(x')\big]$（略去倒像）。取 $\phi_0=0$ 得 $E_1+E_2$（加法）；取 $\phi_0=\frac\pi2$，即光栅沿 $\xi$ 平移四分之一周期 $\frac{\lambda f}{4b}$，得 $E_1-E_2$（减法）。
\end{exam}

\textbf{【适配/变式与数值】}
\begin{itemize}
\item 光栅频率必须与图形间距匹配：$T$ 中的 $\cos(2\pi bu)$ 在频谱面的周期为 $\Lambda=\frac{\lambda f}{b}$。$\lambda=632.8$ nm，$f=200$ mm，$b=5$ mm：$\Lambda=25.3\,\mu$m，减法需平移 $\Lambda/4=6.3\,\mu$m。
\item 间距与光栅不匹配（$b'\neq b$）：$E_1$ 的 $+1$ 级在 $b'-b$，$E_2$ 的 $-1$ 级在 $b-b'$，两者错开 $2(b'-b)$，不能相加。
\item 用 $\sin$ 光栅 $\frac12+\frac12\sin(2\pi bu)$：等于 $\phi_0=-\frac\pi2$，直接是减法（$E_2-E_1$ 或差一个符号，强度相同）。
\item 看到的是强度：$|E_1-E_2|^2$，两图相同处变暗——可用来找两幅图的差异。
\end{itemize}

\begin{formula}{泽尼克相衬与暗场 [F3-s90–92]}\ix{相衬}\ix{泽尼克相衬}\ix{暗场}\ix{弱相位物体}
弱相位物体：$t=\ee^{i\varphi}\approx1+i\varphi$，频谱 $T(u,v)\approx\delta(u,v)+i\Phi(u,v)$（取 $\varphi$ 零均值，直流在频谱面中心）。以下取单位平面波照明、理想 4f 成像，忽略倒像。
\begin{itemize}
\item 明场（完整传递该纯相位场）：$I=|\ee^{i\varphi}|^2=1$，无相位对比。按 [F3-s90] 将场展开到二阶，$-\varphi^2/2$ 与 $i\varphi$ 对强度的二阶贡献抵消。
\item \textbf{暗场}（挡零级）：一阶场 $T'\approx i\Phi$，领先阶强度 $I\approx|\varphi|^2$，看得见但\textbf{非线性} [F3-s91]。
\item \textbf{相衬}（零级乘 $iA$：相移 $90^\circ$、振幅透过 $A$）：$T'=iA\delta+i\Phi$，像面复振幅 $iA+i\varphi=i(A+\varphi)$，
\[
I\approx\boxed{A^2+2A\varphi}\qquad(\text{强度保留到一阶，见 [F3-s92]})
\]
\end{itemize}
相衬强度与相位\textbf{成线性}。相对背景的变化 $\Delta I/I_0\approx2\varphi/A$；线性近似要求 $|\varphi|\ll A$，适当减小 $A$ 可增大相对变化。
（物理说明：$1+i\varphi$ 中直流与散射光相差 $90^\circ$，叠加不产生一阶强度变化；把直流相移 $90^\circ$ 后两者同相或反相，干涉出与 $\varphi$ 成正比的强度。）
\end{formula}

\subsection*{单元：相位物体成像题（明场 / 暗场 / 相衬）}\label{sec:phasecontrast}

\textbf{【题型】}“透明细胞/相位光栅 $t=\ee^{i\varphi(x,y)}$，$|\varphi|\ll1$，在 4f 频谱面分别 (1) 不放东西 (2) 挡零级 (3) 放泽尼克相位板，求像面强度，说明哪种与 $\varphi$ 成线性”。

\textbf{【完整推导】}
\begin{enumerate}
\item 泰勒展开：$\ee^{i\varphi}=1+i\varphi-\frac{\varphi^2}2-\cdots\approx1+i\varphi$（$|\varphi|\ll1$）。
\item 频谱：$\F\{1\}=\delta(u,v)$（只在频谱面中心一点），$\F\{i\varphi\}=i\Phi(u,v)$（分布在外围）。（假设 $\varphi$ 无直流分量；有的话并入常数相位。）
\item 频谱面操作只作用在中心那一点：中心乘 $c$（不放 $c=1$；挡住 $c=0$；相位板 $c=A\ee^{\pm i\pi/2}=\pm iA$）。
\item 像面一阶复振幅（忽略倒像）：$c+i\varphi$ [F3-s90–92]。
\item 按强度所需的阶数计算：明场 $c=1$，完整场给 $I=|\ee^{i\varphi}|^2=1$；暗场 $c=0$，领先阶 $I\approx\varphi^2$；相衬 $c=\pm iA$，保留一阶得 $I\approx A^2\pm2A\varphi$。
\end{enumerate}

\textbf{【适配/数值】}
\begin{itemize}
\item $\varphi_{\max}=0.1$ rad：理想明场 $I=1$；暗场领先阶 $I\lesssim0.01$，且正负 $\varphi$ 一样亮；相衬一阶估计 $A=1$ 时 $I\approx1\pm0.2$（$\pm20\%$），$A=0.5$ 时 $I\approx0.25\pm0.1$（相对背景变化约 $40\%$）。
\item \textbf{正相衬 / 负相衬}：零级相移 $+90^\circ$ 时相位大处更亮（$+2A\varphi$），相移 $-90^\circ$ 时更暗。
\item \textbf{相位不小}（$\varphi\sim1$）：从完整的 $t=\ee^{i\varphi}$ 计算频谱和滤波后的场，再求强度。
\item \textbf{相位光栅} $\varphi=\phi_0\cos\frac{2\pi x}d$：暗场 $I\approx\phi_0^2\cos^2\frac{2\pi x}{d}$，周期变为 $\frac d2$（倍频，同 §\ref{sec:gratingfilter}(b)）；相衬 $I\approx A^2+2A\phi_0\cos\frac{2\pi x}d$，周期 $d$，一阶强度反映相位。
\end{itemize}

\begin{warn}[相衬题易错]
\begin{enumerate}
\item 相位板只盖\textbf{零级}（中心一点），不是盖整个频谱。
\item “相移 $90^\circ$”对应乘 $\ee^{i\pi/2}=i$，不是乘 $\frac\pi2$。
\item 暗场的强度 $\varphi^2$ 与相位不成线性，也分不出 $\varphi$ 的正负。
\end{enumerate}
\end{warn}
